%% kapa: exam=1 round=35 year=2024 period=1 form=A subject=civil_law subject_ko=민법 seq=16 number=16
%% source: Q-Net 감정평가사 기출문제 (cst00309) artlSeq=5213229 file=raw/1008/5213229/1교시 문제지 원본.pdf
%% conversion: kordoc 4.24.51 -> tools/gen1.py (auto); review=proofread-sonnet; flags=none
\begin{question}{16}
취소에 관한 설명으로 옳지 않은 것은? (다툼이 있으면 판례에 따름)
\begin{choices}
  \item 매도인에 의해 매매계약이 적법하게 해제된 후에는 매수인은 그 매매계약을 착오를 이유로 취소할 수 없다.
  \item 법률행위의 취소를 전제로 한 이행거절 가운데는 특별한 사정이 없는 한 취소의 의사표시가 포함된 것으로 볼 수 있다.
  \item 취소할 수 있는 법률행위가 일단 취소된 후에는 취소할 수 있는 법률행위의 추인에 의하여 이를 다시 확정적으로 유효하게 할 수는 없다.
  \item 취소권은 추인할 수 있는 날로부터 3년내에 법률행위를 한 날로부터 10년내에 행사하여야 한다.
  \item 취소할 수 있는 법률행위의 취소권의 행사기간은 제척기간이다.
\end{choices}
\end{question}
\answer{1}
